\section{Policy}\label{sec:policy}

Each component of the IGE points to a different policy, and the data allow two direct tests. Generous unemployment
insurance in the son's teens goes with a smaller stability component, in every sign the data can check, but the
difference is not significant. The search friction faced by the family's group---the job finding of men of its race in
its state---carries a tenth of the IGE, entirely between races. A third implication needs no new design: the
components that grow with the son's age are invisible to evaluations that score children's outcomes around 30.

\subsection{Unemployment insurance in childhood}

If what a layoff passes on is in part the income and the turmoil that follow it, insurance that replaces more of the
lost earnings should weaken the stability component. Table~\ref{tab:ui_stability} measures the generosity of the
state's unemployment insurance as its maximum weekly benefit over its average weekly wage, averaged over the son's ages
12--18 in the state where he lived at 17, and splits sons at the median of generosity relative to their birth cohort.
The IGE is the same in both halves (0.374$^{***}$ and 0.373$^{***}$). The stability component is smaller where
insurance was more generous: 0.035$^{**}$ of the IGE (9.3\%$^{**}$) against 0.050$^{***}$ (13.3\%$^{***}$), and 0.071
against 0.101$^{**}$ at 45--49. The difference, $-4.1$ points of the IGE (s.e.\ 6.0, $p = 0.496$), is the predicted
direction and is not significant.\footnote{Sons born 1958 and later (3,495), whose teens fall in years covered by the
Department of Labor's statutory schedules. Code: \texttt{12b.UIGenerosity.py} (generosity) and
\texttt{82.UIStabilityChannel.py} (\texttt{output/ui\_stability\_channel.md}).}

\begin{table}[!htbp]
\centering
\caption{State unemployment insurance in childhood and the stability channel (PSID)}
\label{tab:ui_stability}
\resizebox{\ifdim\width>\textwidth\textwidth\else\width\fi}{!}{%
\resizebox*{!}{\ifdim\totalheight>0.85\textheight 0.85\textheight\else\totalheight\fi}{%
\begin{tabular}{lccccc}
\toprule
 & \begin{tabular}[b]{@{}c@{}}25--29\end{tabular} & \begin{tabular}[b]{@{}c@{}}35--39\end{tabular} & \begin{tabular}[b]{@{}c@{}}45--49\end{tabular} & \begin{tabular}[b]{@{}c@{}}25--55\end{tabular} & \begin{tabular}[b]{@{}c@{}}Share of\\the IGE,\\25--55 (\%)\end{tabular} \\
 & (1) & (2) & (3) & (4) & (5) \\
\midrule
\multicolumn{6}{l}{\textit{A. Split by UI generosity relative to the son's cohort}} \\
IGE, more generous states & $0.258^{***}$ & $0.433^{***}$ & $0.378^{***}$ & $0.374^{***}$ &  \\
 & (0.036) & (0.050) & (0.077) & (0.035) &  \\
IGE, less generous states & $0.288^{***}$ & $0.381^{***}$ & $0.408^{***}$ & $0.373^{***}$ &  \\
 & (0.038) & (0.056) & (0.088) & (0.044) &  \\
IGE, difference & $-0.031$ & $0.052$ & $-0.031$ & $0.002$ &  \\
 & (0.049) & (0.077) & (0.111) & (0.056) &  \\
Parents' spells (stability), more generous states & $0.037^{**}$ & $0.019$ & $0.071$ & $0.035^{**}$ & $9.3^{**}$ \\
 & (0.016) & (0.022) & (0.043) & (0.016) & (4.4) \\
Parents' spells (stability), less generous states & $0.026^{*}$ & $0.054^{***}$ & $0.101^{**}$ & $0.050^{***}$ & $13.3^{***}$ \\
 & (0.015) & (0.019) & (0.041) & (0.015) & (4.3) \\
Parents' spells (stability), difference & $0.011$ & $-0.035$ & $-0.031$ & $-0.015$ & $-4.1$ \\
 & (0.023) & (0.031) & (0.060) & (0.023) & (6.0) \\
\addlinespace
\multicolumn{6}{l}{\textit{B. Interactions (first column; per standard deviation of each)}} \\
closure share $\times$ UI, ln earnings, sons 25--55 & $0.034^{**}$ &  &  &  &  \\
 & (0.017) &  &  &  &  \\
layoff share $\times$ UI, ln earnings, sons 25--55 & $0.021$ &  &  &  &  \\
 & (0.016) &  &  &  &  \\
mother's unemployment $\times$ UI, son's years unemployed, sons 25--55 & $-0.010^{**}$ &  &  &  &  \\
 & (0.005) &  &  &  &  \\
placebo: father's wage $\times$ UI, ln earnings, sons 25--55 & $-0.004$ &  &  &  &  \\
 & (0.016) &  &  &  &  \\
closure share $\times$ UI, ln earnings, sons 40--55 & $0.057^{**}$ &  &  &  &  \\
 & (0.025) &  &  &  &  \\
layoff share $\times$ UI, ln earnings, sons 40--55 & $0.027$ &  &  &  &  \\
 & (0.030) &  &  &  &  \\
mother's unemployment $\times$ UI, son's years unemployed, sons 40--55 & $-0.009^{*}$ &  &  &  &  \\
 & (0.005) &  &  &  &  \\
placebo: father's wage $\times$ UI, ln earnings, sons 40--55 & $0.000$ &  &  &  &  \\
 & (0.029) &  &  &  &  \\
\bottomrule
\end{tabular}}}
\par\smallskip
\begin{minipage}{\textwidth}\footnotesize \textit{Notes:} UI generosity = the state's maximum weekly benefit over its average weekly wage, mean over the son's ages 12--18 in his state at 17 (DOL Significant Provisions; \texttt{12b.UIGenerosity.py}). Panel A: sons above and below the median of generosity net of five-year cohort means, Table~\ref{tab:passed_on} panel A split in each half, common bootstrap draws. Panel B: the son's outcome on all components (one s.d.), ln parental income, UI and UI $\times$ each element, with state fixed effects, five-year cohort and single-age dummies; standard errors clustered by 1968 family. Sons born 1958 and later (3,495). Source: \texttt{82.UIStabilityChannel.py}; table built by \texttt{90.PaperBRaceNeutralTables.py}. $^{***}$ $p<0.01$, $^{**}$ $p<0.05$, $^{*}$ $p<0.1$.\end{minipage}
\end{table}


The interactions in panel B point the same way. With state fixed effects, a standard deviation of generosity raises the
son's earnings by 0.034$^{**}$ ($p = 0.045$) per standard deviation of the father's closure share, and by 0.057$^{**}$
($p = 0.024$) at 40--55; the layoff interaction is positive and imprecise (0.021, $p = 0.19$); and generosity lowers the
son's share of years unemployed by 0.010$^{**}$ ($p = 0.037$) per standard deviation of the mother's unemployment. The
placebo---generosity interacted with the father's wage, which insurance should not offset---is zero ($-0.004$). Every
sign is that of insurance offsetting the spells component; each estimate alone is weak because generosity varies little
within a state and cohort (a standard deviation of 0.041, against 0.089 in all).\footnote{With the log of the maximum
total benefit instead of the replacement ratio the signs are the same and the estimates weaker (layoff $\times$ UI
0.024, $p = 0.12$). For Black sons the placebo on employment is 0.034$^{**}$ ($p = 0.017$).} A state's schedule is
an imprecise measure of what a given family received; the father's own replacement rate at each layoff would measure it.

\subsection{Group search friction}

A family shares its labor market with other families of its race and place. Table~\ref{tab:group_race} measures the
search friction of the group as the monthly job finding of non-employed men of the son's race in his state during his
teens, from the CPS: 0.193 for White and 0.140 for Black men (difference $-0.054^{***}$). Added to the split of
Section~\ref{sec:passedon}, this measure carries 8.2\%$^{***}$ ($p = 0.004$) of the pooled IGE, against 1.9 percent for
the state's unemployment rate and labor-demand shocks on the same sons. Within race it carries nothing: 0.0 percent of
the White IGE and $-9.1$ percent (s.e.\ 14.5) of the Black.\footnote{Sons born 1976 and later (1,752), whose teens fall
in the CPS monthly files of 1994--2021; the job finding of women of the son's race enters the same component. Annual
re-employment of men of the son's race in the March CPS of 1977--92 carries nothing (0.006, $p = 0.37$). Code:
\texttt{17.AtlasFrictions.py} (measure) and \texttt{84.GroupFrictionByRace.py}
(\texttt{output/group\_friction\_by\_race.md}).} The group component is a contrast between races: Black families have
lower parental income and face lower job finding. Within a race, the friction of the family's place is not passed on.
In these data the group's job finding cannot be told apart from anything else that differs by race, so it says where
policy aimed at groups would act---job-search assistance, anti-discrimination enforcement---not that it would work.

\begin{table}[!htbp]
\centering
\caption{Group search friction: job finding of the son's race in his state (CPS) and the IGE}
\label{tab:group_race}
\resizebox{\ifdim\width>\textwidth\textwidth\else\width\fi}{!}{%
\resizebox*{!}{\ifdim\totalheight>0.85\textheight 0.85\textheight\else\totalheight\fi}{%
\begin{tabular}{lccc}
\toprule
 & \begin{tabular}[b]{@{}c@{}}Pooled\end{tabular} & \begin{tabular}[b]{@{}c@{}}White\end{tabular} & \begin{tabular}[b]{@{}c@{}}Black\end{tabular} \\
 & (1) & (2) & (3) \\
\midrule
\multicolumn{4}{l}{\textit{A. The measure (sons' childhood, 1994--2021)}} \\
Men's monthly job finding, own race, state (White / Black mean) & 0.193 / 0.140 &  &  \\
Black $-$ White & $-0.0537^{***}$ &  &  \\
 & (0.0024) &  &  \\
Gradient in ln parental income &  & $0.0042^{**}$ & $0.0003$ \\
 &  & (0.0021) & (0.0032) \\
\addlinespace
\multicolumn{4}{l}{\textit{B. Share of the earnings IGE carried by the place component}} \\
State unemployment and labour demand (\% of the IGE) & $1.9$ & $1.3$ & $16.1$ \\
 & (1.5) & (1.6) & (10.2) \\
+ job finding of men and women of the son's race (\% of the IGE) & $8.2^{***}$ & $0.0$ & $-9.1$ \\
 & (2.9) & (1.4) & (14.5) \\
\bottomrule
\end{tabular}}}
\par\smallskip
\begin{minipage}{\textwidth}\footnotesize \textit{Notes:} Job finding = monthly probability that a non-employed man (woman) aged 20--45 of the son's race in his state at 17 is employed the next month, CPS basic monthly 1994--2021 matched across months, pooled over the son's ages 12--18 (\texttt{17.AtlasFrictions.py}). Panel A, first row: White / Black means. Panel B: the place component's share of the IGE in Table~\ref{tab:passed_on} panel A's split, on the sons with the measure (born 1976 and later, 1,752). Source: \texttt{84.GroupFrictionByRace.py}; table built by \texttt{90.PaperBRaceNeutralTables.py}. $^{***}$ $p<0.01$, $^{**}$ $p<0.05$, $^{*}$ $p<0.1$.\end{minipage}
\end{table}


\subsection{Transfers}

Years on AFDC or TANF carry 6.6\%$^{**}$ ($p = 0.015$) of the IGE at equal parental income and equal components, and the
share of parental income that came from transfers carries nothing ($-0.001$; Section~\ref{sec:passedon}). Receipt marks
the families whose sons do worse at equal income; the transfer dollars, at equal income, do not show. The split cannot
say whether receipt marks what brought the family onto welfare or what welfare did; it says that the transfer channel
is about which families receive, not about the amount.

\subsection{The age of scoring}

The components that grow with the son's age cannot be seen by evaluations that score outcomes around 30. Long-run
evaluations of childhood policies measure adult earnings in the late twenties or early thirties---Food Stamps
\citep{hoynes2016long}, school finance \citep{jackson2016effects}, the EITC \citep{bastian2018long}, Moving to
Opportunity \citep{chetty2016effects} and the 31--37 window of \citet{chetty2020race}---and welfare calculations project
lifetime earnings from those ages \citep{hendren2020unified}. Between 25--29 and 45--49 the parents' wage component
rises from 0.035 to 0.097$^{***}$ (at 35--39), the father's schooling from 0.014 to 0.075$^{***}$, the spells component
from 0.017$^{**}$ to 0.060$^{***}$, and a teen layoff's cost from nothing to 0.110$^{***}$ a year. A policy that acts on
any of these---schooling, the stability of parents' employment, insurance in the son's teens---would be scored at 30
before most of its effect has appeared. The son's own spells run the other way: a policy that raises young men's
employment shows its effect on the IGE early, and its share falls with age.

\subsection{By race}

Two conditions at the start of the son's career have been proposed for the racial gap in persistence: the labor
market he enters and public employment, where Black sons start more often. Neither creates the racial opening of the
IGE, which is 0.208$^{**}$ among sons who entered tight markets and 0.213$^{***}$ among sons who started in the private
sector (Appendix~\ref{app:entry}). Generous insurance offsets the father's closures most clearly for White sons at
40--55 (0.062$^{**}$, $p = 0.019$) and the mother's unemployment most clearly for Black sons ($-0.020^{**}$ on the son's
years unemployed, $p = 0.011$).
